% eksperymentalne tłumaczenie na włoski

%

\section{Kąty środkowe, wpisane, dopisane} % lang-pl
\section{Angoli al centro, alla circonferenza, tangente-corda} % lang-it


Materiał przytoczony tutaj będzie łatwy do odnalezienia w~dowolnym podręczniku geometrii, dla przykładu u~Guzickiego \cite[s. 11-13, 18, 19]{guzicki_2021}, Audina \cite[s. 74, 75]{audin_2003}, Neugebauera \cite[s. 22, 23]{neugebauer_2018}. % lang-pl
Il materiale riportato qui sarà facile da trovare in un qualunque manuale di geometria, per esempio in Guzicki \cite[pp. 11-13, 18, 19]{guzicki_2021}, Audin \cite[pp. 74, 75]{audin_2003}, Neugebauer \cite[pp. 22, 23]{neugebauer_2018}. % lang-it

\begin{definition}
    Dany jest okrąg $\Gamma$ o~środku w~punkcie $O$, z~którego poprowadzono dwie półproste przecinające okrąg w~punktach $A$ i $B$. % lang-pl
    
    Każdy z dwóch kątów $\angle AOB$ (na ogół jeden z nich jest wklęsły, a drugi wypukły) nazywamy kątem środkowym opartym na łuku $AB$; % lang-pl
    Sia dato un cerchio $\Gamma$ con centro nel punto $O$, dal quale sono tracciate due semirette che intersecano il cerchio nei punti $A$ e $B$. % lang-it
    
    Ciascuno dei due angoli $\angle AOB$ (in generale uno di essi è concavo e l'altro convesso) si chiama angolo al centro insistente sull'arco $AB$; % lang-it
\end{definition}

\begin{definition}
    Dany jest punkt $C$ leżący na okręgu $\Gamma$ oraz dwie półproste o wspólnym początku w punkcie $C$, które przecinają okrąg w punktach $A$ i $B$. % lang-pl
    
    Półproste tworzą jeden kąt wypukły i jeden kąt wklęsły, wewnątrz kąta wypukłego jest zawarty jeden łuk $AB$ okręgu. % lang-pl
    Siano dati un punto $C$ giacente sul cerchio $\Gamma$ e due semirette con origine comune nel punto $C$, che intersecano il cerchio nei punti $A$ e $B$. % lang-it
    
    Kąt wypukły nazywamy kątem wpisanym w okrąg, opartym na łuku $AB$. % lang-pl
    Le semirette formano un angolo convesso e un angolo concavo; all'interno dell'angolo convesso è contenuto un arco $AB$ del cerchio. % lang-it
    
    L'angolo convesso si chiama angolo alla circonferenza, insistente sull'arco $AB$. % lang-it
\end{definition}

\begin{proposition}
\index{kąt!środkowy}% % lang-pl
\index{kąt!wpisany}% % lang-pl
    Kąt środkowy jest dwa razy większy od kąta wpisanego opartego na tym samym łuku. % lang-pl
\index{angolo!al centro}% % lang-it
\index{angolo!alla circonferenza}% % lang-it
    
    L'angolo al centro è il doppio dell'angolo alla circonferenza insistente sullo stesso arco. % lang-it
\begin{figure}[H]\centering%
\begin{comment}
\begin{tikzpicture}[scale=.35]
    \tkzDefPoint(0, 0){Zero}
    \tkzDefPoint(100:5){A}
    \tkzDefPoint(100:3.5){Aa}
    \tkzDefPoint(230:5){B}
    \tkzDefPoint(285:1.5){Oo}
    \tkzDefPoint(340:5){C}
    \tkzDrawCircle[line width=0.5mm](Zero,A)
    \tkzLabelPoint[above left](Zero){$O$}
    \tkzLabelPoint[above](A){$A$}
    \tkzLabelPoint[below left](B){$B$}
    \tkzLabelPoint[below right](C){$C$}
    \tkzDrawPolygons[line width=0.3mm](A,B,Zero,C)
    \tkzMarkAngle(B,Zero,C)
    \tkzMarkAngle(B,A,C)
    \tkzDrawPoints[size=3,color=black,fill=red!50](A,B,C,Zero)
    \tkzLabelPoint[anchor=center](Aa){$\alpha$}
    \tkzLabelPoint[anchor=center](Oo){$2\alpha$}
\end{tikzpicture}
\end{comment}
    \caption{kąt środkowy jest dwa razy większy od wpisanego} % lang-pl
    
    \caption{l'angolo al centro è il doppio dell'angolo alla circonferenza} % lang-it
\end{figure}
\end{proposition}
% PRZECZYTANO: https://en.wikipedia.org/wiki/Inscribed_angle

To będzie (III.20). % lang-pl

Questo sarà (III.20). % lang-it

\begin{corollary}
    \label{maybe_3_28}
    Kąty wpisane oparte na tym samym łuku (ogólniej: na równych łukach) są równe. % lang-pl
    
    Gli angoli alla circonferenza insistenti sullo stesso arco (più in generale: su archi uguali) sono uguali. % lang-it
\end{corollary}

Neuegebauer (a poza nim żaden autor podręcznika do geometrii) nazwie nie wiedzieć czemu powyższy wniosek twierdzeniem Apoloniusza. % lang-pl
Neuegebauer (e, oltre a lui, nessun autore di manuali di geometria) chiamerà, chissà perché, il corollario precedente teorema di Apollonio. % lang-it
\begin{corollary} % lang-it
    L'angolo insistente su una semicirconferenza è retto. % lang-it
\end{corollary} % lang-it


\begin{corollary} % lang-pl
    Kąty oparty na półokręgu jest prosty. % lang-pl
\end{corollary} % lang-pl
Grecki pisarz Diogenes Laertios przypisze to stwierdzenie Talesowi, chociaż nazwa przyjmie się jedynie w krajach anglosaskich. % lang-pl
Lo scrittore greco Diogene Laerzio attribuirà questa affermazione a Talete, benché il nome si diffonderà soltanto nei paesi anglosassoni. % lang-it
\index[persons]{Tales z Miletu}%
\index[persons]{Laertios, Diogenes}%
Pojawia się u Pompego \cite[s. 37, 38]{pompe_2022}. % lang-pl
Compare in Pompe \cite[pp. 37, 38]{pompe_2022}. % lang-it


Można wykorzystać je do konstrukcji stycznej. % lang-pl
\index{styczna}% % lang-pl
Si può usare questo fatto per costruire una tangente. % lang-it
\index{tangente}% % lang-it

\begin{corollary}
    \label{ab_twice_pi}
    Kąty wpisane oparte na dwóch różnych łukach $AB$ dają w sumie kąt półpełny. % lang-pl
    
    Gli angoli alla circonferenza insistenti sui due diversi archi $AB$ danno in somma un angolo piatto. % lang-it
\end{corollary}

Wykorzystamy ten wniosek później (fakt \ref{prp_incircle}) do opisania, na których czworokątach można opisać okrąg. % lang-pl

Useremo questo corollario più avanti (fatto \ref{prp_incircle}) per descrivere su quali quadrilateri si può circoscrivere un cerchio. % lang-it

\begin{proposition}
    Kąt dopisany, między styczną do okręgu a cięciwą przechodzącą przez punkt styczności, jest równy kątowi wpisanemu w ten okrąg, opartemu na odpowiednim łuku: $\angle ABC = \angle CAD$. % lang-pl
    
    \index{kąt!dopisany} % lang-pl
    L'angolo tangente-corda, compreso tra la tangente a un cerchio e la corda passante per il punto di tangenza, è uguale all'angolo alla circonferenza in quel cerchio, insistente sull'arco corrispondente: $\angle ABC = \angle CAD$. % lang-it
    \index{angolo!tangente-corda} % lang-it
\begin{figure}[H] \centering
\begin{comment}
\begin{tikzpicture}[scale=.4]
    \tkzDefPoint(0, 0){Zero}
    \tkzDefPoint(50:5){A}
    \tkzDefPoint(230:5){B}
    \tkzDefPoint(300:5){C}
    \tkzDefLine[tangent at=A](Zero) \tkzGetPoint{DD}
    \tkzDefPointsBy[projection=onto A--DD](C){D}
    \tkzDrawCircle[line width=0.5mm](Zero,A)
    \tkzDrawLines[add= 1 and 1, line width=0.3mm](A,D)
    \tkzDrawPolygons[line width=0.3mm](A,B,C)
    \tkzMarkAngle[arc=l,size=1.2,mark=||](C,B,A)
    \tkzMarkAngle[arc=l,size=1.2,mark=||](C,A,D)
    \tkzDrawPoints[size=3,color=black,fill=black!50](A,B,C,D)
    \tkzLabelPoint[above left](Zero){$O$}
    \tkzLabelPoint[above right](A){$A$}
    \tkzLabelPoint[below left](B){$B$}
    \tkzLabelPoint[below right](C){$C$}
    \tkzLabelPoint[above right](D){$D$}
\end{tikzpicture}
\end{comment}
    \caption{kąt dopisane jest równy wpisanemu} % lang-pl
    
    \caption{l'angolo tangente-corda è uguale all'angolo alla circonferenza} % lang-it
\end{figure}
\end{proposition}

\todofoot{https://www.deltami.edu.pl/2015/09/styczna-i-cieciwa/}

%